% introduction.tex -- BNR-2026-012

\section{Introduction}
\label{sec:introduction}

Data availability is a foundational requirement for any Layer~2 (L2)
architecture that posts validity proofs rather than full transaction data
to the settlement layer. In a validium or zkEVM design, the prover
demonstrates correct state transitions via zero-knowledge proofs, but the
underlying transaction data must remain retrievable so that any party can
independently reconstruct the L2 state. A Data Availability Committee
(DAC) provides this guarantee by distributing batch data among a set of
trusted nodes that attest to its availability before the proof is
submitted on-chain~\cite{starkware2023da, albassam2018fraud}.

Enterprise deployments introduce constraints absent from public L2
systems. First, batch data contains commercially sensitive information
(supply-chain records, financial transactions, industrial telemetry) that
must not be exposed to individual committee members. Second, the
committee is operated by a small consortium of enterprise-grade nodes
($n \leq 10$), not a large validator set. Third, recovery latency and
storage overhead must be minimized to support real-time traceability
workflows with sub-second finality requirements.

Our prior work (RU-V6~\cite{ruv6}) addressed enterprise privacy by
applying Shamir $(2,3)$-secret sharing directly to batch data. While
this achieves information-theoretic privacy, it incurs $3.87\times$
storage overhead and $O(k^2)$ recovery cost per field element---scaling
poorly as committee size increases to the 5--7 nodes required for
production fault tolerance.

This paper presents a hybrid DAC architecture that restructures the
privacy and redundancy layers:

\begin{enumerate}
  \item \textbf{AES-256-GCM encryption} provides computational data
    privacy (254-bit key entropy).
  \item \textbf{Reed-Solomon $(5,7)$ erasure coding} distributes the
    ciphertext into $n = 7$ chunks, any $k = 5$ of which suffice for
    reconstruction, achieving $1.40\times$ storage overhead.
  \item \textbf{Shamir $(5,7)$-secret sharing} protects only the 32-byte
    AES key, preserving information-theoretic key secrecy at negligible
    cost.
\end{enumerate}

The combination achieves computational data privacy (no single node or
coalition of fewer than $k$ nodes can decrypt the batch) with
information-theoretic key secrecy (fewer than $k$ Shamir shares reveal
zero information about the AES key). This layered approach yields
$36\times$ faster attestation, $2{,}613\times$ faster recovery, and
$2.77\times$ lower storage overhead compared to RU-V6.

\subsection{Contributions}

The contributions of this work are:

\begin{itemize}
  \item A hybrid AES+RS+Shamir DAC architecture that decouples data
    privacy from data redundancy, achieving both with minimal overhead.
  \item A Go implementation using the SIMD-optimized
    \texttt{klauspost/reedsolomon} library~\cite{klauspost}, validated
    against production systems (MinIO, Storj, CockroachDB).
  \item Comprehensive benchmarks across five batch sizes (10\,KB--5\,MB)
    with 50+ replications per configuration, including attestation
    latency, storage overhead, recovery time, failure tolerance,
    availability probability, and privacy validation.
  \item Benchmark reconciliation with published results from EigenDA
    V2~\cite{eigenda2024v2}, Arbitrum AnyTrust~\cite{anytrust},
    Celestia~\cite{celestia}, and Avail~\cite{avail}.
  \item Six formal invariants (INV-DA-P1 through INV-DA-P6) for
    downstream TLA+ specification and Coq verification.
\end{itemize}

\subsection{Paper Organization}

Section~\ref{sec:related} surveys production DAC architectures and
erasure coding techniques. Section~\ref{sec:system-model} defines the
system model, threat model, and invariants.
Section~\ref{sec:methodology} describes the hybrid protocol and
experimental setup. Section~\ref{sec:results} presents benchmark
results across all metrics. Section~\ref{sec:discussion} analyzes
trade-offs, limitations, and production readiness.
Section~\ref{sec:conclusion} concludes with recommendations for the
Basis Network zkEVM L2 pipeline.
